\chapter{Assemblers}

Onderstaande tabel geeft een overzicht van assemblers waarvan bekend is dat ze met de MEGA65 werken.
Voor algemeen gebruik raden we {\bf ACME} aan, omdat het goede ondersteuning biedt voor de
45GS02-instructieset, open source is en bovendien in C is geschreven. Dat laatste
betekent dat de assembler op termijn naar de MEGA65 geporteerd zou kunnen worden om er native te draaien.

\begin{longtable}{ | l | l | l | l |}\hline
Naam     & 45GS02 & Bron & Verwijzing \\\hline
ACME     &  ja    & C    & \url{https://sourceforge.net/projects/acme-crossass}\\
KickAss  &  ja    & Java & \href{https://gitlab.com/jespergravgaard/kickassembler65ce02}{gitlab.com/jespergravgaard/kickassembler65ce02}\\
LLVM-MOS &  ja    & C++  & \url{https://llvm-mos.org/wiki/Welcome}\\
Ophis    &  ja    & Python & \url{https://github.com/michaelcmartin/Ophis}\\
BSA      &  ja    & C    & \url{https://github.com/Edilbert/BSA}\\
CA65     &  nee\footnote{Onze fork van CA65 (onderdeel van CC65) herkent de CPU van de MEGA65 correct, maar heeft geen expliciete ondersteuning voor de mogelijkheden van de processor} & C & \url{https://github.com/mega65/cc65}\\\hline
\end{longtable}

De {\bf BSA}-assembler wordt momenteel gebruikt om de {\bf MEGA65.ROM} te bouwen.
Het grootste deel van de broncode is geschreven in de syntaxis van de oude
{\bf BSO}-assembler (Boston Systems Office), die in de jaren 1989--1991 werd
gebruikt door softwareontwikkelaars aan de C65. De {\bf BSA}-assembler heeft een
compatibiliteitsmodus waarmee deze oude broncodes met weinig of geen aanpassingen
kunnen worden geassembleerd. De {\bf BSA}-assembler bevat momenteel alleen een
beschrijving van de commando's in de C-broncode van de assembler.

Het {\bf LLVM-MOS}-project wordt geleverd met een GNU-compatibele macro-assembler
en -disassembler die alle mnemonieken van de 45GS02 ondersteunt. Het kan als
zelfstandig hulpprogramma worden gebruikt met het commando {\bf llvm-mc}.
